\begin{table}[H]
\centering
\small
\setlength{\tabcolsep}{4pt}
\begin{tabular}{llcrrr}
\toprule
Model & Task & Total cells & $P=8$ & $P=16$ & $P=32$ \\
\midrule
LR & Coarse & 3,200 & 0.969 [0.958, 0.976] & 0.972 [0.963, 0.978] & 0.975 [0.971, 0.980] \\
LR & Coarse & 6,400 & 0.967 [0.950, 0.975] & 0.973 [0.967, 0.977] & 0.975 [0.972, 0.979] \\
LR & Coarse & 12,800 & 0.965 [0.948, 0.974] & 0.973 [0.961, 0.978] & 0.978 [0.975, 0.981] \\
LR & Fine & 3,200 & 0.566 [0.551, 0.583] & 0.574 [0.560, 0.594] & 0.579 [0.567, 0.593] \\
LR & Fine & 6,400 & 0.588 [0.568, 0.602] & 0.594 [0.579, 0.613] & 0.599 [0.581, 0.612] \\
LR & Fine & 12,800 & 0.604 [0.584, 0.615] & 0.611 [0.592, 0.626] & 0.615 [0.602, 0.624] \\
\midrule
\bottomrule
\end{tabular}
\caption{Fixed-total training-cell scaling. Each row holds the total number of sampled training cells constant while reallocating them across independent patients. Entries are subject-balanced pooled, patient-disjoint OOF Macro-F1 means with empirical 95\% intervals over 20 matched subset seeds. For each score, every held-out subject's confusion matrix is normalized to unit mass before pooling; this is distinct from both cell-weighted pooling and mean within-subject Macro-F1. All outer-test cells are evaluated.}
\label{tab:fixed_total_scaling}
\end{table}
